\bookchapter{Capitolo 55: Progetto 2 — l'app del meteo (versione in rete)}

\chapterintro{Il registro delle spese è «ginnastica interna», l'app del meteo è «aprirsi al mondo» — deve trattare con internet. In questo progetto imparerai: richiedere un'API gratuita, fare il parsing del JSON, gestire i tre stati dell'interfaccia «caricamento / successo / errore» e come nascondere bene la chiave API. Completandolo avrai la capacità di realizzare qualunque «app connessa» — notizie, tassi di cambio, traduzioni: lo schema è sempre lo stesso.}

\booksection{Richiedere un'API meteo gratuita}

Per esercitarsi da soli, ti consiglio OpenWeatherMap oppure QWeather (il «meteo del vento armonioso», molto usato in Cina) — entrambe hanno una quota gratuita, basta registrarsi. Prendiamo OpenWeatherMap come esempio:

\begin{enumerate}
  \item Registrati nella pagina API keys genera una Key (apikey);
  \item Il piano gratuito può chiamare l'interfaccia Current Weather (meteo attuale): api.openweathermap.org/data/2.5/weather? q=beijing\&appid=LA\_TUA\_KEY\&units=metric;
  \item Prima provala una volta nel browser, per vedere com'è fatto il JSON restituito — è la «preparazione» da fare sempre prima di collegare qualunque API.
\end{enumerate}

\textbf{Attenzione alle trappole: la chiave non va mai scritta a mano nel codice}\par

La API Key è come la «password del tuo portafoglio»: se la scrivi nel codice e lo mandi su GitHub, entro poche ore i crawler la scandagliano e la usano al posto tuo. La strada giusta: mettila in local.properties (un file già ignorato da gitignore) e leggila nel codice attraverso il buildConfig di Gradle. Cerca «best practice per nascondere le API Key su Android» e segui la guida.

\booksection{Il livello di rete: Retrofit + parsing del JSON}

Prima definisci il «modello dati» sul calco del JSON restituito dall'API; i campi che non corrispondono si ignorano senza pietà:

\begin{lstlisting}[language=Kotlin]
// JSON: {"weather":[{"description":"Sereno"}],"main":{"temp":26.5,"humidity":40},"name":"Pechino"}
data class WeatherResponse(
    val weather: List<Weather>,
    val main: Main,
    val name: String
)
data class Weather(val description: String)
data class Main(val temp: Double, val humidity: Int)

interface WeatherApi {
    @GET("data/2.5/weather")
    suspend fun getWeather(
        @Query("q") city: String,
        @Query("appid") key: String,
        @Query("units") metric: String = "metric",
        @Query("lang") lang: String = "it"
    ): WeatherResponse
}
\end{lstlisting}

Gson infila da solo il JSON dentro queste data class — se i nomi dei campi coincidono, la magia avviene da sola.

\booksection{I tre stati dell'interfaccia: caricamento / successo / errore}

L'anima di un'app connessa è la «gestione degli stati». Rete lenta, rete assente, nome della città sbagliato... in ogni situazione l'interfaccia deve rendere conto all'utente. Si rappresenta con eleganza con una sealed class (classe sigillata):

\begin{lstlisting}[language=Kotlin]
sealed class UiState {
    object Loading : UiState()                    // spinner in rotazione
    data class Success(val data: WeatherResponse) : UiState()  // successo
    data class Error(val message: String) : UiState()          // errore
}

// Nel ViewModel:
fun search(city: String) {
    viewModelScope.launch {
        _state.value = UiState.Loading
        _state.value = try {
            UiState.Success(api.getWeather(city))
        } catch (e: Exception) {
            UiState.Error("Città non trovata — provare con un altro nome?")
        }
    }
}

// Nell'interfaccia:
when (val s = state) {
    is UiState.Loading -> CircularProgressIndicator()   // spinner
    is UiState.Success -> WeatherContent(s.data)        // mostra
    is UiState.Error -> Text(s.message)                 // errore
}
\end{lstlisting}

Questo schema a tre stadi Loading/Success/Error è il «modello tuttofare» delle interfacce connesse: prima imparalo a memoria, poi per ogni futura app connessa non ti resta che seguirne il calco.

\booksection{Il codice d'interfaccia sul campo: disegnare il meteo}

Con gli stati pronti, l'interfaccia si presenta così (frammenti chiave):

\begin{lstlisting}[language=Kotlin]
@Composable
fun WeatherContent(data: WeatherResponse) {
    Column(
        Modifier.fillMaxSize().padding(24.dp),
        horizontalAlignment = Alignment.CenterHorizontally
    ) {
        Text(data.name, fontSize = 28.sp, fontWeight = FontWeight.Bold)
        Text(
            "${data.main.temp}°C",
            fontSize = 72.sp,
            fontWeight = FontWeight.Light
        )
        Text(
            data.weather.firstOrNull()?.description ?: "Sconosciuto",
            fontSize = 20.sp,
            color = Color.Gray
        )
        Spacer(Modifier.height(16.dp))
        Row {
            Text(" Umidità ${data.main.humidity}%")
            Spacer(Modifier.width(24.dp))
            Text(" Percepiti ${data.main.temp}°C")
        }
    }
}
\end{lstlisting}

Qualche dettaglio da imparare: firstOrNull() prende in sicurezza il primo elemento (l'elenco potrebbe essere vuoto), ?: “sconosciuto” dà il testo di ripiego — il principio della null safety vale anche nell'interfaccia. La temperatura usa un corpo grande come 72.sp per fare figura, e il peso «leggero» (Light) rende il numero elegante: è la ricetta classica delle app meteo.

\booksection{Dettagli di interfaccia e d'interazione}

\booksection{Dettagli di interfaccia e d'interazione}

\begin{itemize}
  \item Campo di ricerca: un TextField in alto, digiti il nome della città e tocchi cerca (imposta l'azione «Cerca» sulla tastiera, Invio per cercare);
  \item Scheda meteo: temperatura a corpo grande (26,5°), descrizione del tempo (sereno), umidità, nome della città;
  \item Città preferite: salva in DataStore le città che guardi più spesso e alla prossima apertura le cerca da solo — così ripassi di passo il capitolo 6;
  \item Tira per aggiornare: in Compose usa PullToRefreshBox (offerto da Material 3) e la sensazione al tocco diventa subito «professionale».
\end{itemize}

\begin{tcbinfo}{Consiglio: le immagini danno vita all'interfaccia}
Abbina a ogni tipo di meteo un emoji o un'icona (), mappandoli con when(description) — tre righe di codice e il calore umano dell'interfaccia raddoppia. Puoi anche caricare una libreria di icone meteo di terze parti, ma nella fase di esercizio gli emoji bastano e avanzano.

\end{tcbinfo}

\begin{tcbpractice}{Mani in pasta: la checklist di collaudo dell'app del meteo}
(1) La ricerca del meteo funziona; (2) stacca la rete: l'app segnala l'errore con eleganza senza bloccarsi (prova con la modalità aereo!); (3) digiti una città che non esiste e il messaggio resta gentile; (4) metti due città nei preferiti e, riavviata l'app, ci sono ancora; (5) ruotando lo schermo lo stato non si perde. Solo superando tutti e cinque i punti il progetto si può dire promosso.

\end{tcbpractice}

\begin{tcbsummary}{Riassunto del capitolo}
\begin{itemize}
  \item Prima prova l'API nel browser e impara a leggere il JSON, poi scrivi il codice: «prepararsi in anticipo» non è mai tempo perso
  \item Nascondi la chiave API in local.properties e non farla mai finire nel repository
  \item Interfaccia Retrofit + data class: se i campi corrispondono, il JSON si analizza da solo
  \item Lo schema a tre stati Caricamento/Successo/Errore è il coltellino svizzero delle schermate di rete
  \item Rete assente, input vuoto, città inesistente: la gestione degli errori è il voto minimo per un'app connessa
\end{itemize}

\end{tcbsummary}
